\section{The operator square}
\label{sec:operators}

We now fix the parameters of a square supplied by
Proposition~\ref{prop:square-locality}.  The aim is to describe the functions
at one cell that are visible both from the rest of its row and from the rest
of its column.  We shall prove that this common subspace is spanned by
unimodular functions satisfying multiplicative relations in both directions.
The proof first treats three groups of rows and three groups of columns;
varying the groups will then recover the individual coordinates.
Operator representations of joinings have a substantial precedent in
Ryzhikov's intertwining work~\cite{Ryzhikov1994,Ryzhikov1997}; the
particular rectangle identities and exact-line conclusion below are
proved from the conditional kernels of Section~\ref{sec:arrays}.

\subsection{Active spaces and exact relations}

\begin{definition}\label{def:active-space}
Let $\lambda$ be a probability law on a finite product
$S_1\times\cdots\times S_m$ of standard spaces, with coordinate variables
$Y_1,\ldots,Y_m$ and marginals $\mu_1,\ldots,\mu_m$.
The \emph{active space} at coordinate $i$ is the closed subspace
\[
 \mathcal A_i(\lambda)
 =\overline{\Span}\bigl\{
   \E_\lambda[F(Y_j:j\ne i)\mid Y_i]:F\text{ bounded and measurable}
 \bigr\}\subseteq L^2(\mu_i).
\]
An \emph{exact tuple} is a tuple of measurable functions
$\chi_i:S_i\to\T$ for which
\begin{equation}\label{eq:exact-tuple}
 \prod_{i=1}^m\chi_i(Y_i)=c\qquad\lambda\text{-almost surely},
 \qquad c\in\T.
\end{equation}
An \emph{exact line} at coordinate $i$ is a one-dimensional subspace
$\C\chi_i$ occurring in an exact tuple.  Representatives differing by a
constant of modulus one define the same line.
\end{definition}

The active spaces contain constants and are invariant under complex
conjugation.  Their orthogonal projections therefore commute with
conjugation.  Every exact line lies in the active space: solve
\eqref{eq:exact-tuple} for its representative and condition on the chosen
coordinate.

\begin{lemma}\label{lem:exact-lines}
Suppose $m\ge3$ and every $m-1$ coordinates of $\lambda$ are independent.
Distinct exact lines at any coordinate are orthogonal.  An exact line at
one coordinate has a unique tuple of partner lines at the other coordinates.
The exact tuples of lines form an abelian group under coordinatewise
multiplication, and projection to any one coordinate is a group
isomorphism onto its group of exact lines.  Each of these groups is
countable when the marginal $L^2$ spaces are separable.  Every nonconstant
exact line has zero mean.
\end{lemma}

\begin{proof}
For an exact tuple put $a_i=\abs{\int\chi_i\dd\mu_i}$.  Solving the product
relation for $\chi_i$ and using independence of the remaining coordinates
gives
\[
 a_i=\prod_{j\ne i}a_j,\qquad 0\le a_i\le1.
\]
If some $a_i=0$, at least one other $a_j$ is zero, and the same equations
then force all $a_j$ to be zero.  Otherwise put $b_i=-\log a_i$.
The equations give $2b_i=\sum_j b_j$ for every $i$; since $m\ge3$, all
$b_i$ vanish.  Thus either all means have modulus one, in which case all
the functions are constant almost surely, or all means are zero.

The quotient of two exact tuples is again exact.  If their lines differ
at one coordinate, its quotient function is nonconstant, so all quotient
means vanish.  In particular the two representatives are orthogonal at
that coordinate.  If their lines agree at one coordinate, the quotient
there is constant, forcing constant quotients everywhere.  This proves
uniqueness of partners.  Products and conjugates preserve exact tuples;
uniqueness makes multiplication of lines well defined and makes each
coordinate projection bijective.  A family of nonzero mutually orthogonal
lines in a separable Hilbert space is countable.  Comparing any exact tuple
with the constant tuple gives the final assertion.
\end{proof}

For the rest of the section, fix a parameter square $E$ in the conull set
of Proposition~\ref{prop:square-locality}, and suppress $E$ from the
notation.  Partition the rows into three nonempty groups $P,A,B$, and the
columns into three nonempty groups $r,s,t$.  A grouped cell $Y_{i,l}$ is
the tuple of states in the corresponding rectangle.  Write $\mu_{i,l}$
for its marginal and $H_{i,l}=L^2(\mu_{i,l})$.
The square locality statement supplies the following facts, for every
ordering of the three row groups and of the three column groups:
\begin{enumerate}
\item Each grouped row and each grouped column has a three-coordinate
law whose coordinates are pairwise independent.
\item Given the whole grouped row $P$, any two grouped columns are
independent.  The conditional law of each uses only its own entry in $P$.
\item Given the whole row $P$ and all columns except $t$, the conditional
law of $Y_{A,t}$ is its own row-law kernel given $Y_{A,r},Y_{A,s}$.
\end{enumerate}
These statements concern permutations of groups within each axis; no
interchangeability of the two axes is assumed.
Let $V_{i,l}$ and $W_{i,l}$ be the active spaces for the grouped row and
grouped column laws at $(i,l)$, with projections $P^V_{i,l}$ and
$P^W_{i,l}$, respectively.

\subsection{Conditional Markov operators}

Set
\[
 z=Y_{P,r},\qquad c=Y_{P,s},\qquad d=Y_{P,t},\qquad
 x=Y_{A,r},\qquad y=Y_{B,r}.
\]
The triples $(z,c,d)$ and $(z,x,y)$ are pairwise independent.  Moreover,
$(c,d)$ and $(x,y)$ are conditionally independent given $z$: the law of
column $r$ given the whole row $P$ depends only on $z$.
It follows that each individual variable in either pair is independent
of the whole opposite pair.  For example, for bounded $f,g$,
\[
 \E[f(x)g(c,d)]
 =\E\bigl[\E[f(x)\mid z]\E[g(c,d)\mid z]\bigr]
 =\E[f(x)]\E[g(c,d)].
\]
In particular, each triple formed from one variable of a pair and the
whole opposite pair has product law.

Define four conditional Markov operators, with the arrow indicating the
direction in which functions are mapped:
\begin{equation}\label{eq:square-operators}
\begin{aligned}
 L_A(x)&:H_{A,t}\longrightarrow H_{A,s},&
 L_B(y)&:H_{B,t}\longrightarrow H_{B,s},\\
 J_s(c)&:H_{B,s}\longrightarrow H_{A,s},&
 J_t(d)&:H_{B,t}\longrightarrow H_{A,t}.
\end{aligned}
\end{equation}
For instance,
\[
 (L_A(x)h)(Y_{A,s})
 =\E[h(Y_{A,t})\mid Y_{A,r}=x,Y_{A,s}]
\]
in the grouped row law.  Pairwise independence makes the source and
target marginals independent of the conditioning value, so these are
contractions between the fixed Hilbert spaces displayed above.
Their adjoints are the reverse conditional Markov operators.  Kernels
make their matrix coefficients measurable; separability then gives
strong measurability on each vector.

Figure~\ref{fig:operator-square} fixes the four cell Hilbert spaces and
the five conditioned cells. The first identity in
Lemma~\ref{lem:rectangle} compares the two paths from the lower-right
vertex to the upper-left. The second compares the paths from the
lower-left to the upper-right, using the reverse row operators.
\begin{figure}[htbp]
  \centering
  \begin{tikzpicture}[
    conditioned cell/.style={
      draw=black!45, dashed, fill=black!7,
      minimum width=25mm, minimum height=10mm,
      inner sep=3pt
    },
    operator cell/.style={
      draw=black, fill=white, line width=.6pt,
      minimum width=25mm, minimum height=10mm,
      inner sep=3pt
    },
    operator arrow/.style={-{Stealth[length=2.2mm]}, line width=.7pt},
    operator label/.style={fill=white, inner sep=2pt, font=\small}
  ]
    \node at (0,1) {$r$};
    \node at (4,1) {$s$};
    \node at (8,1) {$t$};

    \node[anchor=east] at (-1.8,0) {$P$};
    \node[anchor=east] at (-1.8,-2) {$A$};
    \node[anchor=east] at (-1.8,-4) {$B$};

    \node[conditioned cell] (Pr) at (0,0) {$z=Y_{P,r}$};
    \node[conditioned cell] (Ps) at (4,0) {$c=Y_{P,s}$};
    \node[conditioned cell] (Pt) at (8,0) {$d=Y_{P,t}$};
    \node[conditioned cell] (Ar) at (0,-2) {$x=Y_{A,r}$};
    \node[conditioned cell] (Br) at (0,-4) {$y=Y_{B,r}$};

    \node[operator cell] (As) at (4,-2) {$H_{A,s}$};
    \node[operator cell] (At) at (8,-2) {$H_{A,t}$};
    \node[operator cell] (Bs) at (4,-4) {$H_{B,s}$};
    \node[operator cell] (Bt) at (8,-4) {$H_{B,t}$};

    \draw[operator arrow] (At.west) --
      node[operator label,above] {$L_A(x)$} (As.east);
    \draw[operator arrow] (Bt.west) --
      node[operator label,above] {$L_B(y)$} (Bs.east);
    \draw[operator arrow] (Bs.north) --
      node[operator label,left] {$J_s(c)$} (As.south);
    \draw[operator arrow] (Bt.north) --
      node[operator label,right] {$J_t(d)$} (At.south);
  \end{tikzpicture}
  \caption{The grouped operator square at a fixed parameter $E$.
    The five shaded cells record the conditioned values $z,c,d,x,y$.
    Each arrow acts on functions in its source Hilbert space by
    conditional expectation in the indicated row or column law.
    Reversing an arrow represents its adjoint.}
  \label{fig:operator-square}
\end{figure}


\begin{lemma}[The rectangle identities]\label{lem:rectangle}
For the grouped square just described, almost surely in $(z,c,d,x,y)$,
\begin{equation}\label{eq:rectangle-identities}
 L_A(x)J_t(d)=J_s(c)L_B(y),\qquad
 L_A(x)^*J_s(c)=J_t(d)L_B(y)^*.
\end{equation}
The identities also hold in every ordering of the row and column groups.
\end{lemma}

\begin{proof}
Condition on $z,c,d,x,y$.  The marginal law of the two remaining entries
in column $s$ is the column law given $c$, and similarly column $t$ has
its law given $d$.  Indeed, given row $P$, each of these columns is
independent of column $r$.  Denote the corresponding two joint $L^2$
spaces by $\mathcal H_s$ and $\mathcal H_t$.
Let $Q:\mathcal H_t\to\mathcal H_s$ be the conditional Markov operator
between them under this five-cell conditioning, and let
$\iota_{A,s},\iota_{B,s},\iota_{A,t},\iota_{B,t}$ be the isometric
embeddings of the individual cell spaces into the appropriate joint
spaces.

The own-row kernel assertion gives
\[
 Q\iota_{A,t}=\iota_{A,s}L_A,\qquad
 Q\iota_{B,t}=\iota_{B,s}L_B.
\]
The same assertion with $s$ and $t$ reversed gives
\[
 Q^*\iota_{A,s}=\iota_{A,t}L_A^*,\qquad
 Q^*\iota_{B,s}=\iota_{B,t}L_B^*.
\]
Also $J_s=\iota_{A,s}^*\iota_{B,s}$ and
$J_t=\iota_{A,t}^*\iota_{B,t}$.  Consequently
\[
 L_AJ_t=\iota_{A,s}^*Q\iota_{B,t}=J_sL_B,
 \qquad
 L_A^*J_s=\iota_{A,t}^*Q^*\iota_{B,s}=J_tL_B^*.
\]
Countable dense sets of tests make these operator equalities simultaneous
on a conull set.  The grouped locality assertions hold for all the
indicated orderings, so the same proof applies to each of them.
\end{proof}

We next identify the subspaces on which the rectangle identities carry
information.  For a measurable operator family, its \emph{essential
closed range span} means the smallest closed subspace containing its
ranges for almost every parameter value.

\begin{lemma}\label{lem:active-reduction}
In every grouped cell, the projections $P^V$ and $P^W$ commute.  Every row
Markov operator intertwines the column-active projections at its source
and target; every column Markov operator intertwines the row-active
projections.  Put
\[
 D_{i,l}=V_{i,l}\cap W_{i,l},\qquad
 P^D_{i,l}=P^V_{i,l}P^W_{i,l}.
\]
Each $D_{i,l}$ is invariant under conjugation, and $P^D_{i,l}$ commutes
with conjugation.
All incident row and column operators intertwine the $P^D$ projections.
For either choice of source in a row or column, the restrictions to $D$
have essential closed range span equal to the target $D$; the same holds
for the reverse operators.
\end{lemma}

\begin{proof}
First, the essential closed range span of $L_A(x)$ is $V_{A,s}$.
Conditional expectations of product tests onto $Y_{A,s}$ have the form
\begin{equation}\label{eq:weighted-ranges}
 \int w(x)L_A(x)h\dd\mu_{A,r}(x),
\end{equation}
where $w,h$ are bounded.  Such tests span the active space by density of
product functions in the $L^2$ space of the other two coordinates.
A vector orthogonal to every integral in \eqref{eq:weighted-ranges} is
orthogonal to $L_A(x)h$ almost everywhere, by varying $w$.  A countable
dense set of $h$ then makes it orthogonal to the whole range for almost
every $x$.  The converse is immediate.  This proves the assertion, and
the same argument applies to every incident operator in either direction.
In particular each such operator has its range in its own-direction
active space and annihilates the orthogonal complement of the source
active space.

By the first rectangle identity, $J_s(c)L_B(y)$ has range in $V_{A,s}$.
The pair $(c,y)$ is independent.  Fubini and the essential range span of
$L_B$ therefore give
$J_s(c)V_{B,s}\subseteq V_{A,s}$ for almost every $c$.
The adjoint of the second identity is
$J_s(c)^*L_A(x)=L_B(y)J_t(d)^*$.
Independence of $(c,x)$ gives the reverse inclusion for the adjoint.
Together these say
\begin{equation}\label{eq:column-intertwining}
 P^V_{A,s}J_s(c)=J_s(c)P^V_{B,s}.
\end{equation}
Likewise, the first rectangle identity and independence of $(x,d)$
show that $L_A(x)W_{A,t}\subseteq W_{A,s}$; the second identity and
independence of $(x,c)$ give the adjoint assertion.  Thus
\begin{equation}\label{eq:row-intertwining}
 P^W_{A,s}L_A(x)=L_A(x)P^W_{A,t}.
\end{equation}
The same reasoning works at every edge.

Equation~\eqref{eq:column-intertwining} shows that $P^V_{A,s}$ preserves
the essential range span of $J_s$, which is $W_{A,s}$.  As an orthogonal
projection it also preserves $W_{A,s}^{\perp}$, and therefore commutes
with $P^W_{A,s}$.  This proves commutation in every cell.  Combining
\eqref{eq:row-intertwining} with the own-row active-space reduction gives
$P^D_{A,s}L_A=L_AP^D_{A,t}$, and similarly for all other edges.
Finally, projecting the essential range span $V_{A,s}$ onto $D_{A,s}$
shows that the restricted ranges of $L_A$ span $D_{A,s}$:
$P^D_{A,s}\Ran L_A=\Ran(L_AP^D_{A,t})$.
The corresponding argument handles each edge and its adjoint.
\end{proof}

We have obtained commuting projections and a rectangle of contractions
between their intersections.  The next two lemmas show that the parameter
dependence of this rectangle is carried by an abelian group of unitaries.
That description will produce exact functions at the corner cell.

\subsection{Fixed positive parts and an abelian unitary group}

Restrict the four operators in \eqref{eq:square-operators} to the $D$
spaces.  Temporarily denote them, in the same order, by $F(x)$, $G(y)$,
$K(c)$ and $N(d)$.  Their equations are
\begin{equation}\label{eq:restricted-rectangle}
 FN=KG,\qquad F^*K=NG^*.
\end{equation}
All four families and their adjoints have the essential range spans
proved in Lemma~\ref{lem:active-reduction}.

\begin{lemma}[Fixed squares]\label{lem:fixed-squares}
For the four restricted families in \eqref{eq:restricted-rectangle},
each source square and target square is independent of its parameter
almost surely, and each is injective on its $D$ space.  The four polar
parts are measurable families of unitaries between the $D$ spaces and
satisfy the first rectangle identity.
\end{lemma}

\begin{proof}
The two rectangle equations and their adjoints imply
\begin{align}
 KK^*F&=FNN^*,& K^*KG&=GN^*N,\label{eq:squares-first}\\
 F^*FN&=NG^*G,& FF^*K&=KGG^*.
 \label{eq:squares-second}
\end{align}
The parameter triples in these four equations are, respectively,
$(c,x,d)$, $(c,y,d)$, $(x,d,y)$ and $(x,c,y)$.  Each has product law by
the independence observations preceding Lemma~\ref{lem:rectangle}.

For the first equation, compare two independent values $c,c'$ using the
same $x,d$.  Fubini gives
\[
 \bigl(K(c)K(c)^*-K(c')K(c')^*\bigr)F(x)=0
\]
for almost every $(c,c',x)$.  The essential range span of $F$ is its whole
target $D$ space, so the two target squares agree.  The adjoint equation
and the essential range span of $F^*$ similarly make $NN^*$ fixed.
The second equation in \eqref{eq:squares-first} makes $K^*K$ and $N^*N$
fixed.  The two equations in \eqref{eq:squares-second}, using the range
spans of $N,K$ and their adjoints, make $F^*F,G^*G,FF^*,GG^*$ fixed.
These comparisons are operator comparisons: countable dense vector tests
first give a common conull set, after which boundedness gives equality
on every vector.

If $v$ is in the kernel of the fixed square $F^*F$, then $F(x)v=0$ for
almost every $x$.  Hence $v$ is orthogonal to the ranges of almost every
$F(x)^*$, whose essential span is the whole source $D$.  Thus $v=0$.
The same reasoning proves injectivity of all source and target squares.
Each operator consequently has zero kernel and dense range, so its polar
part has full initial and final spaces and is unitary.

To preserve the rectangle identity, put $A=|F|$ and $B=|G|$.  From
$F^*FN=NG^*G$ we obtain $A^2N=NB^2$.  Polynomial approximation to the
square-root function on $[0,1]$ gives $AN=NB$.  Write $F=U_FA$ and
$G=U_GB$.  The equation $FN=KG$ becomes
\[
 (U_FN-KU_G)B=0.
\]
The range of the injective self-adjoint $B$ is dense, so the bounded
operator in parentheses vanishes.  If $N=U_N|N|$ and $K=U_K|K|$, then
\[
 U_FN=(U_FU_N)|N|,\qquad
 KU_G=(U_KU_G)(U_G^*|K|U_G)
\]
are polar decompositions of the same operator.  Uniqueness yields
$U_FU_N=U_KU_G$.  This argument uses no inverse bound for the positive
parts.  Finally, polar parts are strongly measurable on vectors, for
example by the bounded approximations
$F(F^*F+\varepsilon\id)^{-1/2}$ as $\varepsilon\downarrow0$.
\end{proof}

Denote these polar parts by $a(x),b(y),c_0(c),d_0(d)$, respectively.  They
satisfy
\begin{equation}\label{eq:polar-rectangle}
 a(x)d_0(d)=c_0(c)b(y).
\end{equation}
The subspaces are nonzero, since constants belong to every $D$ space.

\begin{lemma}[Normalization of the polar rectangle]\label{lem:polar-group}
There are fixed unitary identifications of the four $D$ spaces with one
separable Hilbert space $\mathcal H$ such that the four polar variables
in \eqref{eq:polar-rectangle} take values in one closed abelian subgroup
$\mathcal G$ of the unitary group of $\mathcal H$.  Their common support
is $\mathcal G$.  The variables $a,b$ have a common law $\rho$, invariant
under left multiplication by every element of $\mathcal G$, and there
is a measurable $\mathcal G$-valued function $h(z)$ such that
\begin{equation}\label{eq:corner-ratio}
 a(x)b(y)^{-1}=c_0(c)d_0(d)^{-1}=h(z).
\end{equation}
The group $\mathcal G$ has a simultaneous orthonormal eigenbasis.
Monomials in its diagonal circle coordinates, evaluated at $h(z)$,
span $L^2(\sigma(h(z)))$ and are exact functions for both the row and
column three-cell laws at $(P,r)$.
\end{lemma}

\begin{proof}
Give the unitary spaces the strong operator topology.  On unitaries,
inversion and multiplication are continuous, and these spaces are
separable metrizable.  Conditional on a typical fixed value $z_0$, the
pair $(a,b)$ and the pair $(c_0,d_0)$ are independent, while each individual
marginal equals its unconditional marginal.  The continuous identity
\eqref{eq:polar-rectangle} therefore holds throughout the product of
the two conditional pair supports.
Choose one point $(a_*,b_*)$ from the first support and one point
$(c_*,d_*)$ from the second.  They satisfy $a_*d_*=c_*b_*$.  Identify the
four vertex spaces by these reference unitaries, starting at the lower
right vertex.  The two possible identifications of the upper left vertex
agree by this equation, so all four reference edge values become identity.
Use these fixed identifications from now on.

The conditional pair supports at $z_0$ now contain $(\id,\id)$.
Fixing the column pair at this point in \eqref{eq:polar-rectangle}
forces $a=b$ throughout the other pair support.  Fixing the row pair
similarly forces $c_0=d_0$.  The two row marginals therefore have a common
law $\rho$ and the two column marginals a common law $\tau$; these are
also their unconditional laws.  The supports of $\rho$ and $\tau$ both
contain identity.  Substitution of diagonal pair values into
\eqref{eq:polar-rectangle} shows that these two supports commute
elementwise.

For general $z$, conditional independence of the two pairs and their
fixed individual marginals imply
\begin{equation}\label{eq:translated-laws}
 \rho=u\rho v^{-1}
 \quad\text{for almost every actual pair }(u,v)=(c_0,d_0).
\end{equation}
Indeed, after fixing $z,c_0,d_0$, the equation says $a=c_0bd_0^{-1}$,
and $a,b$ each have law $\rho$.  Unconditionally, $c_0,d_0$ are independent,
so their pair law is $\tau\otimes\tau$.  The map
$(u,v)\mapsto u\rho v^{-1}$ is continuous for weak convergence of
probability measures: apply strong-topology continuity and dominated
convergence to each bounded continuous test.  Thus
\eqref{eq:translated-laws} holds for every $u,v\in\supp\tau$.
The symmetric argument gives $\tau=u\tau v^{-1}$ for every
$u,v\in\supp\rho$.

Taking $v=\id$ shows that each support is preserved by left translation
by elements of the other.  Because identity belongs to both, they contain
one another.  Write their common closed support as $\mathcal G$.
For $g\in\mathcal G$ we have $g\mathcal G=\mathcal G$.  This proves closure
under products, and $\id\in g\mathcal G$ supplies $g^{-1}\in\mathcal G$.
Thus $\mathcal G$ is a group.  It is abelian by the earlier commutation,
and \eqref{eq:translated-laws} gives left invariance of $\rho$.
Equation~\eqref{eq:polar-rectangle} now gives equality of the two ratios
in \eqref{eq:corner-ratio}.  Given $z$, the ratios are independent and
equal almost surely, hence have a common point-mass conditional law.
A countable determining family in the standard unitary group makes its
point measurable in $z$, giving $h$.

Choose a positive injective trace-class operator $Q$ on $\mathcal H$ and
form the trace-norm integral
\[
 Q_0=\int_{\mathcal G}UQU^*\dd\rho(U).
\]
The integrand is trace-norm continuous under strong-unitary convergence,
as follows first for finite-rank $Q$ and then by trace-norm approximation.
The integral is positive and trace-class, and left invariance of $\rho$
makes it commute with every element of $\mathcal G$.  For $v\ne0$ each
quadratic form $\ip{v}{UQU^*v}$ is positive; hence so is
$\ip{v}{Q_0v}$.  Thus $Q_0$ is injective.  Its finite-dimensional positive
eigenspaces span $\mathcal H$.  Each is invariant under the commuting
unitaries, so simultaneous diagonalization in each gives a full
orthonormal eigenbasis $(e_n)$.

Write $Ue_n=\gamma_n(U)e_n$, with $\gamma_n(U)\in\T$.
These coordinates are continuous characters.  On diagonal unitaries,
coordinatewise convergence is equivalent to strong convergence, by
approximation of any vector by its finite basis truncations.  The
coordinates therefore generate the Borel information of $\mathcal G$.
Trigonometric polynomials on every finite coordinate torus are dense in
$L^2$ of its actual probability law.  Increasing conditional expectations
onto the finite-coordinate fields then show that the functions
\[
 \gamma(h(z)),\qquad
 \gamma=\prod_{n\in J}\gamma_n^{m_n},\quad
 J\text{ finite},\quad m_n\in\Z,
\]
span $L^2(\sigma(h(z)))$.  No identification of the law of $h$ as Haar
measure is needed for this density statement.
Finally, \eqref{eq:corner-ratio} gives
\[
 \gamma(h(z))
 =\gamma(a(x))\overline{\gamma(b(y))}
 =\gamma(c_0(c))\overline{\gamma(d_0(d))}.
\]
Each factor is unimodular.  Multiplying by the conjugate factors shows
that $\gamma(h(z))$ is exact both in column $r$ and in row $P$.
\end{proof}

\subsection{Completeness of the exact functions}

Lemma~\ref{lem:polar-group} gives exact functions at the corner.  To show
that they span the entire common active subspace, we must also prove
that no other dependence remains in $D_{P,r}$.  The following projection
identity supplies this reverse inclusion.

\begin{lemma}[Projection of tensor contractions]\label{lem:contraction-completeness}
In any three-cell row of the grouped square, abbreviate the cell spaces
by $H_r,H_s,H_t$ and their common-active projections by $P_r,P_s,P_t$.
For $g\in H_s$, $h\in H_t$, set
\[
 C_r(g,h)=\E[g(Y_s)h(Y_t)\mid Y_r].
\]
Then
\begin{equation}\label{eq:contraction-projections}
 P_rC_r(g,h)=C_r(g,P_th)=C_r(P_sg,h)=C_r(P_sg,P_th).
\end{equation}
The contractions with both inputs in their $D$ spaces lie in $D_r$ and
span it.  The same assertions hold for each grouped column.
Moreover, with $h(z)$ from Lemma~\ref{lem:polar-group},
\begin{equation}\label{eq:corner-completeness}
 D_{P,r}=L^2(\sigma(h(z))).
\end{equation}
Consequently the intersection at every grouped cell is spanned by lines
exact simultaneously for its three-cell row and column laws.
\end{lemma}

\begin{proof}
Since $(Y_s,Y_t)$ has product marginal, the bilinear contraction is
well defined on $H_s\times H_t$ and satisfies
\begin{equation}\label{eq:contraction-bound}
 \norm{C_r(g,h)}_2\le\norm{g}_2\norm{h}_2.
\end{equation}
For bounded inputs let $L_{r\leftarrow t}(v)$ be the row Markov operator
from $t$ to $r$ given $Y_s=v$.  Pairwise independence gives the
vector-valued integral formula
\[
 C_r(g,h)=\int g(v)L_{r\leftarrow t}(v)h\dd\mu_s(v).
\]
Lemma~\ref{lem:active-reduction} says
$P_rL_{r\leftarrow t}(v)=L_{r\leftarrow t}(v)P_t$.
Moving the projection through the integral proves
$P_rC_r(g,h)=C_r(g,P_th)$.
Using instead the edge from $s$ to $r$ conditioned on $Y_t$ proves
$P_rC_r(g,h)=C_r(P_sg,h)$.
The bound \eqref{eq:contraction-bound} extends both identities to all
$L^2$ inputs, including unbounded projected tests.  Idempotence then gives
\[
 P_rC_r(g,h)=P_rC_r(g,P_th)=C_r(P_sg,P_th),
\]
proving \eqref{eq:contraction-projections}.  Restricted contractions are
fixed by $P_r$ and hence lie in $D_r$.  Unrestricted product contractions
span the row-active space; projecting that space onto $D_r$ proves the
restricted spanning assertion.  The column proof is identical.

Apply this result first in row $A$.  For $g\in D_{A,s}$ and
$h\in D_{A,t}$, its contraction onto $x=Y_{A,r}$ is
\[
 C_{A,r}(g,h)(x)=\ip{\overline g}{L_A(x)h}.
\]
Conjugation preserves $D_{A,s}$.  On the restricted spaces,
$L_A(x)=a(x)A_0$ for a fixed positive operator $A_0$ by
Lemma~\ref{lem:fixed-squares}.  Thus every displayed coefficient is a
function of $a(x)$ alone.  The space $L^2(\sigma(a(x)))$ is closed,
being the range of a conditional-expectation projection, so the spanning
assertion gives
\[
 D_{A,r}\subseteq L^2(\sigma(a(x))),\qquad
 D_{B,r}\subseteq L^2(\sigma(b(y))).
\]
Apply the column version of the spanning assertion in column $r$.
Given $z$, we have $a=h(z)b$ and $b$ has its fixed marginal law $\rho$.
For bounded measurable $f,g$ the contraction of $f(a)$ and $g(b)$ onto
$z$ is therefore
\[
 \int_{\mathcal G} f(h(z)u)g(u)\dd\rho(u),
\]
a function of $h(z)$.  The product marginal of $(x,y)$, the bound
\eqref{eq:contraction-bound}, and bounded approximation extend this
conclusion to the two $D$ spaces and their contraction span.  Closure
therefore yields $D_{P,r}\subseteq L^2(\sigma(h(z)))$.
In the other direction, the monomials of
Lemma~\ref{lem:polar-group} are exact for both laws, hence lie in
$D_{P,r}$, and they span $L^2(\sigma(h(z)))$.
This proves \eqref{eq:corner-completeness}.  We have used no assumption
that $D$ was already closed under multiplication.  Permuting row groups
and column groups places any grouped cell at the corner and proves the
last assertion.
\end{proof}

\subsection{From grouped cells to individual coordinates}

Return to the individual square, with row tuples $p_a$, column tuples
$q_j$, and marginal $\eta_{p_a[j]}$ at cell $(a,j)$.  Let $V_{a,j}$ be
the active space there for the full row law $\lambda_{p_a}$ and
$W_{a,j}$ the active space for the full column law $\lambda_{q_j}$.

\begin{proposition}[Common exact lines]\label{prop:common-exact-lines}
For almost every parameter square $E$, the following assertions hold
simultaneously at all individual cells.
\begin{enumerate}
\item The orthogonal projections onto $V_{a,j}$ and $W_{a,j}$ commute.
Their common range $V_{a,j}\cap W_{a,j}$ is spanned by the lines exact
for both the full $k$-coordinate row law $\lambda_{p_a}$ and the full
$k$-coordinate column law $\lambda_{q_j}$.
\item For distinct columns $s,t$, let $L^a_{s\leftarrow t}(v)$ be the
row Markov operator from cell $(a,t)$ to $(a,s)$, conditioned on the
remaining $k-2$ row entries $v$.  Then, for almost every $v$,
\[
 P^W_{a,s}L^a_{s\leftarrow t}(v)
 =L^a_{s\leftarrow t}(v)P^W_{a,t}.
\]
Column operators likewise intertwine the corresponding row-active
projections.
\item If $(\chi_j)_{j\in I}$ is an exact tuple for row $a$, with
$\prod_j\chi_j(Y_{a,j})=c$, then
\begin{equation}\label{eq:exact-partner-transport}
 L^a_{s\leftarrow t}(v)\chi_t
 =c\!\prod_{j\ne s,t}\overline{\chi_j(v_j)}\,
   \overline{\chi_s}.
\end{equation}
In particular, membership of one partner line in the column-active
space is equivalent to membership of any other partner line in its
column-active space.  The analogous assertion holds with rows and
columns exchanged.
\end{enumerate}
\end{proposition}

\begin{proof}
Isolate row $a$ and column $j$ as singleton groups, and split each of
their complements into two nonempty groups.  At the isolated cell, the
row-active and column-active spaces do not depend on the chosen splits:
the other two grouped cells still contain exactly all the other entries
of that row or column.  Lemma~\ref{lem:active-reduction} gives commutation,
and Lemma~\ref{lem:contraction-completeness} gives a spanning family of
simultaneous exact lines for the two split three-cell laws. We take this
family to contain all such lines: each is in the intersection, and
orthogonality leaves no additional line outside the spanning family.

Fix the row-group partition and vary the column-group partition.  Every
resulting family is a subset of the mutually orthogonal exact lines of
the same fixed split column law, by Lemma~\ref{lem:exact-lines}.  All have
the same closed span $V_{a,j}\cap W_{a,j}$.  Two subsets of one orthogonal
family with the same closed span must coincide: a line belonging only to
one would be orthogonal to the other span.  Now fix the initial column
grouping and vary the row grouping.  The same argument, using the fixed
split row law, shows that the initial family is exact for every split
in both directions.

For $k=3$ the grouped laws already have individual coordinates.  For
$k>3$, fix a unimodular representative $\chi_j$ of one of the lines and
consider its row completions.  For each split $S$ of $I\setminus\{j\}$,
there are unimodular functions $a_S,b_S$ and $c_S\in\T$ with
\[
 \chi_j(Y_{a,j})a_S(Y_{a,S})b_S(Y_{a,I\setminus(S\cup\{j\})})=c_S.
\]
Absorb $c_S^{-1}$ into $a_S$.  All the resulting completion functions
then equal $\overline{\chi_j(Y_{a,j})}$ under the row law, so they define
the same function $\phi$ of the other coordinates.  Those $k-1$
coordinates have product marginal.  For each singleton split, $\phi$
lies in a one-dimensional line in that coordinate tensored with the
remaining $L^2$ spaces.  The orthogonal projections to these individual
line factors commute and all fix $\phi$.  Their product has the full
one-dimensional tensor product as range, so $\phi$ is a pure product
of its unimodular individual factors, up to one scalar.  This proves
exactness for the full row law.  The column argument is the same.
Conversely, every line exact for both full laws is active for both and
belongs to their intersection.  This proves part~(1).

For part~(2), make the chosen row a singleton group, the chosen columns
$s,t$ singleton groups, and put all remaining columns in the third
group.  Split the remaining rows into two nonempty groups.  The
intertwining in Lemma~\ref{lem:active-reduction} is exactly the stated
identity for this grouping.  To obtain the column assertion, isolate
the two chosen rows instead and apply the already proved column
intertwining; there is again no interchange of axes.

Finally, solving the exact product relation for $\chi_t$ and conditioning
gives \eqref{eq:exact-partner-transport}.  Its scalar has modulus one.
The projections $P^W$ commute with conjugation.  Thus the intertwining
in part~(2) carries membership of $\C\chi_t$ in $W_{a,t}$ to membership
of $\C\chi_s$ in $W_{a,s}$.  The reverse row Markov operator gives the
converse.  The same proof applies to column partners.
\end{proof}

All constructions in this section were performed at a fixed typical
parameter square.  Kernel and operator identities can be imposed using
countable determining tests, and there are only finitely many groupings
and orientations.  The resulting conull set of parameter squares can
therefore be chosen measurable before any selection of exact lines.
The unitary reference choices and diagonal bases were needed only for
the pointwise proof: the conclusions of
Proposition~\ref{prop:common-exact-lines} concern the original conditional
laws and do not require those choices to vary measurably with $E$.
In the next section these intrinsic conclusions, together with the
parameter geometry, will identify the exact lines on individual fibres.
